\scanpage{6}
\mspage{-3.1-}

\uline{Chapter} \uline{3}. \uline{An} \uline{elementary} \uline{exemple}

The purpose of this chapter is to give, independently of the forgoing theory, an exhaustive analysis of a single special class of actions. This has two objectives: first of all, the concepts of pre-stable and stable are worked out in a representative non-trivial case. Secondly, by attacking this case directly, we can circumvent the difficulties involved in extending the previous work to semi-simple groups in characteristic $p$ (cf. Preface). And, in fact, we obtain the result over the ring of integers, hence, \uline{a} \uline{fortiori}, over any field. This will enable us later to construct the so-called arithmetic schemes of moduli of abelian varieties.

We have no use for a ground field in this chapter: all the schemes that we will consider will be of finite type over $\operatorname{Spec}(\handwritten{\mathbb Z})$. $\handwritten{\mathbb A}^{n}$, $\handwritten{\mathbb P}_n$, and $\mathrm{PGL}(n)$ stand for the usual schemes \uline{over} \uline{$\handwritten{\mathbb Z}$}. The action we shall study is the canonical action:

\[
\sigma_m:\mathrm{PGL}(n)\times(\handwritten{\mathbb P}_n)^{m+1}\handwritten{\longrightarrow}(\handwritten{\mathbb P}_n)^{m+1}
\]

\margin{\handwritten{\unsure{--}}}

i.e. of projective transformations on sequences of points in projective space.

\handwritten{\fbox}{\S1.} The first step is to obtain a large number of\handwritten{\unsure{$\mid$}}invariant affine open subsets in $(\handwritten{\mathbb P}_n)^{m+1}$. To this end, let $p_0,\ldots,p_m$ denote the projections of $(\handwritten{\mathbb P}_n)^{m+1}$ onto its $m+1$ factors. Let $L_i=p_i^*(\underline{\mathrm{o}}_{\handwritten{\mathbb P}}(1))$. Now unfortunately, $\underline{\mathrm{o}}_{\handwritten{\mathbb P}}(1)$ admits no $\mathrm{PGL}(n)$-linearization with respect to the action of $\mathrm{PGL}(n)$ on $\handwritten{\mathbb P}_n$. However, if

% archive sheet number: 5
